\section{3. Estimands and extraction
models}\label{estimands-and-extraction-models}

\subsection{3.1 Primary class: fitted gap-intercept
models}\label{primary-class-fitted-gap-intercept-models}

The primary analysis includes only models that return an intercept
parameter conventionally denoted \(E_g\).

The fractional-power family is

\[
\alpha(E)=\frac{A(E-E_g)^p}{E}, \qquad E>E_g,
\]

with either a fitted exponent or fixed \(p=0.5\), \(0.7\), the
source-panel exponent, or \(1.0\).

The historical Chu Kane-region model \citep{Chu1994} is

\[
\alpha(E)=\alpha_g\exp\left[\sqrt{\beta(x,T)(E-E_g)}\right],
\]

where \(\beta(x,T)\) is fixed by the published source relation and the
candidate is used only inside its declared composition and temperature
domain.

These models do not necessarily estimate an identical microscopic
quantity. Their intercept spread is therefore described as
\textbf{fitted-model dependence}, not as repeated-measurement noise.

\textbf{Table 2. Provenance and interpretation of fitted-intercept
candidates.}

\begin{longtable}[]{@{}
  >{\raggedright\arraybackslash}p{(\columnwidth - 6\tabcolsep) * \real{0.2500}}
  >{\raggedright\arraybackslash}p{(\columnwidth - 6\tabcolsep) * \real{0.2500}}
  >{\raggedright\arraybackslash}p{(\columnwidth - 6\tabcolsep) * \real{0.2500}}
  >{\raggedright\arraybackslash}p{(\columnwidth - 6\tabcolsep) * \real{0.2500}}@{}}
\toprule\noalign{}
\begin{minipage}[b]{\linewidth}\raggedright
Candidate
\end{minipage} & \begin{minipage}[b]{\linewidth}\raggedright
Fitted and fixed quantities
\end{minipage} & \begin{minipage}[b]{\linewidth}\raggedright
Provenance or rationale
\end{minipage} & \begin{minipage}[b]{\linewidth}\raggedright
Interpretation and restriction
\end{minipage} \\
\midrule\noalign{}
\endhead
\bottomrule\noalign{}
\endlastfoot
Free exponent & \(A\), \(E_0\), and \(p\) fitted; \(0.2\leq p\leq1.5\) &
Fractional-power family used in the source analysis, with the exponent
released & Flexible fitted intercept; practical fit-window conditioning
must be reported \\
\(p=0.5\) & \(A\) and \(E_0\) fitted & Predeclared fixed-exponent
sensitivity candidate & Not asserted as the correct microscopic
transition law \\
\(p=0.7\) & \(A\) and \(E_0\) fitted & Predeclared exponent within the
source-reported range & Sensitivity candidate only \\
Source-panel \(p\) & \(A\) and \(E_0\) fitted; \(p=0.872\) or \(0.849\)
fixed & Exponent recorded from the corresponding source-panel analysis &
Provenance-bound fitted intercept \\
\(p=1\) & \(A\) and \(E_0\) fitted & Predeclared fixed-exponent
sensitivity candidate & Not asserted as the correct microscopic
transition law \\
Chu Kane-region & \(\alpha_g\) and \(E_0\) fitted; \(\beta(x,T)\) fixed
& Chu \emph{et al.} (1994), used within \(0.170\leq x\leq0.443\) and
\(77\leq T\leq300\) K & Historical observation model; boundary hits are
not treated as identified \\
\end{longtable}

The ensemble was fixed before the robustness analyses. No candidate was
added or removed because of its fitted intercept.

\subsection{3.2 Secondary class: fixed-absorption operational
coordinates}\label{secondary-class-fixed-absorption-operational-coordinates}

The first interpolated crossings at selected values of \(\alpha\) are
retained as operational coordinates. A coordinate such as

\[
E_{\alpha=1000}
\]

can be useful for process control or detector modeling when the
definition is held fixed. It is not automatically an estimator of the
latent zone-center gap. Fixed-\(\alpha\) coordinates are therefore
analyzed and tabulated separately and are excluded from the primary
fitted-model span and from empirical-equation ranking.
